% ── Section 3: Dataset ────────────────────────────────────────────────────────
\section{Dataset and Preprocessing}
\label{sec:dataset}

% ── 3.1 Dataset ───────────────────────────────────────────────────────────────
\subsection{Dataset}
\label{sec:dataset:source}

We use \texttt{cluster-trace-gpu-v2026}, a production GPU cluster trace
released by Alibaba alongside the OSDI~2026 paper~\cite{alibaba2026trace}.
The trace captures Alibaba Serverless Infrastructure (ASI), a hyperscale
scheduling platform that co-locates training, inference, and batch workloads.
Four fact tables are provided: \texttt{server\_hourly} (per-server GPU
inventory at hourly resolution), \texttt{job\_execution\_summary} (one row
per pod lifecycle with scheduling and priority metadata),
\texttt{network\_hourly} (per-server bandwidth for a 7-day window), and
\texttt{pod\_hourly} (fine-grained per-pod utilisation, distributed
separately as a 351\,GB archive).
The trace spans 185~days and peaks at 155{,}410~GPUs across 37{,}707~servers,
making it the largest publicly available GPU cluster trace to date.
Alibaba anonymises the data in two ways consequential to our methodology:
absolute calendar dates are suppressed in favour of a monotonic \texttt{day}
counter, and server region labels are omitted; we address both through
sensitivity analyses in Section~\ref{sec:dataset:limitations}.

% ── 3.2 Preprocessing ─────────────────────────────────────────────────────────
\subsection{Preprocessing}
\label{sec:dataset:preprocessing}

\paragraph{Temporal indexing.}
The \texttt{hour} column in \texttt{server\_hourly} is stored as a
zero-padded \texttt{VARCHAR} (e.g., \texttt{"07"}).
We cast it to \texttt{INTEGER} and derive a monotonic relative-hour index
$\mathit{rel\_hour} = \mathit{day} \times 24 + \mathit{hour\_int}$,
spanning $[0,\,4439]$ across the 185-day trace window.

\paragraph{Server tier classification.}
We partition physical servers into three capacity tiers by \texttt{gpu\_count}:
\emph{small} ($\leq 2$~GPUs, typically single-card inference hosts),
\emph{medium} ($3$--$8$~GPUs, standard training nodes),
and \emph{large} ($>8$~GPUs, high-memory multi-node pods).
This taxonomy captures the distinct scheduling regimes observed across
workload types in the trace.

\paragraph{Deferral eligibility.}
A pod is marked \texttt{is\_deferrable} when three conditions hold jointly:
its \texttt{job\_type\_public} is not \texttt{online\_inference} (no
user-facing latency requirement), its \texttt{priority\_class} is LP
(low priority), and its \texttt{ready\_status} is \texttt{TRUE} (the pod
was fully resourced but not yet dispatched).
These are the pods CADS (Section~\ref{sec:cads}) is permitted to shift.

\paragraph{Truncation handling.}
Pods first observed in state \textsc{running} at $\mathit{rel\_hour}=0$
are flagged left-truncated (true start time unknown); pods still
\textsc{running} at $\mathit{rel\_hour}=4439$ are flagged right-truncated.
Both sets are excluded from duration and fragmentation-gap calculations
but retained for aggregate GPU-hour accounting.

\paragraph{Thermal design power assignment.}
We assign TDP from public Nvidia and AMD spec sheets for all known GPU
models (A100~=~400\,W, H100~=~700\,W, H20~=~500\,W, L20~=~275\,W, etc.).
Alibaba's proprietary XPU accelerators (XPU-A through XPU-E) carry no
published TDP; we assign conservative estimates of 250--350\,W based on
comparable datacenter accelerator classes and report a $\pm30\%$
sensitivity sweep across all XPU rows in
Section~\ref{sec:carbon_analysis}.

% ── 3.3 Limitations ───────────────────────────────────────────────────────────
\subsection{Limitations}
\label{sec:dataset:limitations}

\paragraph{Suppressed calendar dates.}
Absolute dates are replaced by a relative counter, preventing direct
alignment with historical grid carbon intensity signals.
We run all carbon analyses under three calendar anchors
(1~Jan, 1~Apr, and 1~Jul~2025) and report the resulting range of
\mbox{CO\textsubscript{2}}-equivalent estimates throughout.

\paragraph{Absent region labels.}
Server locations are not disclosed, precluding assignment to specific
electricity grids.
We instead model carbon intensity under four synthetic Chinese grid
profiles (North, East, South, Northwest) parameterised by published
annual averages and diurnal load curves, spanning the plausible range
from coal-heavy to hydro-mixed generation mixes
(Section~\ref{sec:carbon_analysis}).

\paragraph{Unavailable pod-hourly utilisation.}
The \texttt{pod\_hourly} table was unavailable during our analysis;
intra-hour resource utilisation is therefore unobservable.
Our fragmentation estimates derive from \texttt{server\_hourly}
allocation records, so the reported allocation gap is a
\emph{lower bound} on true idle GPU time.
We expect Phase~3b analysis on \texttt{pod\_hourly} to reveal
additional utilisation-level waste not captured here.

\paragraph{XPU thermal uncertainty.}
The unknown TDP of Alibaba's XPU family is the largest single
uncertainty in our energy estimates.
The $\pm30\%$ sensitivity sweep in Section~\ref{sec:carbon_analysis}
bounds the resulting figures and demonstrates that qualitative
findings are stable across the full range of plausible XPU power
envelopes.